\section{Norm cancellation for odd multiples of the source point}
\label{sec:odd-norm-cancellation}

Let $K=\mathbb Q(\sqrt2)$ and retain the source curve
$E:y^2=x^3-(11-6\sqrt2)x$ and point
\[
 P=\left(57/8-7\sqrt2/2,-165/8+417\sqrt2/32\right).
\]
The odd-product identity adds heights of abscissas inside $K$.
Passing from that product to its rational norm need not preserve
the height. The failure occurs already at an odd source multiple,
and its additive size is unbounded along odd multiples.

\begin{proposition}\label{cc:odd-norm-cancellation}
Write $N=\operatorname{Norm}_{K/\mathbb Q}$ and let $h$ be the
absolute logarithmic Weil height. Then
\[
        h(Nx([5]P))=2h(x([5]P))-2\log137.
\]
For every integer $r\ge0$ put $n_r=5\cdot137^r$. Then
\[
         2h(x([n_r]P))-h(Nx([n_r]P))
                         \ge2(r+1)\log137.
\]
In particular, the defect is unbounded even when every index is odd.
\end{proposition}
\begin{proof}
The prime $137$ splits in $K$, with $\sqrt2\equiv31$ and
$\sqrt2\equiv-31$ at its two places. The exact certificate
\texttt{odd-norm-cancellation-verification.json} computes $[5]P$
by both $P+[4]P$ and $[2]P+[3]P$. It verifies the curve equation
and, using Hensel roots modulo $137^8$, gives its two abscissa
valuations as $2$ and $-2$.

For the exact first identity, write $x([5]P)=(a+b\sqrt2)/c$ with
primitive integers $a,b,c$ and $c>0$. Put
\[
 g=\gcd(c^2,a^2-2b^2),\qquad
 g_0=\gcd(c^2,2ac,a^2-2b^2).
\]
The certificate records all three integers and checks $g/g_0=137^2$.
Both real conjugates of $x([5]P)$ exceed one; the same certificate
verifies this by the rational inequalities
$a/c>1$ and $(a/c-1)^2>2(b/c)^2$.
The primitive minimal polynomial has leading coefficient $c^2/g_0$,
whereas the reduced denominator of its rational norm is $c^2/g$.
The archimedean contributions agree after multiplying the degree-two
height by two. The difference of the two heights is therefore
$\log(g/g_0)=2\log137$.

For the infinite assertion, put $T_0=(0,0)$ and $Q=[5]P$.
At the place where $v(x(Q))=-2$, the point $Q$ is formal of
depth one. At the other place, $Q+T_0$ is formal of depth one,
because $x(Q+T_0)=-(11-6\sqrt2)/x(Q)$ and the coefficient is a
unit at $137$. The formal logarithm is an isometry on this
unramified neighborhood. Multiplication by $137^r$ increases
the depth by exactly $r$. Since $137^r$ is odd,
\[
          [137^r](Q+T_0)=[137^r]Q+T_0.
\]
The two valuations of $x([n_r]P)$ are therefore respectively
$-2(r+1)$ and $2(r+1)$. Their contribution to the rational norm
valuation cancels, whereas their contribution to twice the
height over $K$ is $2(r+1)\log137$. At every other place the
norm-height inequality has nonnegative defect. Summing proves
the asserted lower bound.
\end{proof}

This excludes an exact norm-height identity and a uniform additive
error bound on the odd cyclic coset. It does not exclude an
effective polynomial comparison between the two heights. No such
comparison, and no deletion of a detector on this basis, is proved.
